\chapter{Discussion, extended}
\label{ch:discussion}

\noindent\emph{\Cref{ch:discussionmain} is the argument. This appendix adds
the detail that did not fit there.}

\section{The clearance model itself is not yet trustworthy}

The most uncomfortable single finding of the hardware programme belongs
here rather than in the results. Asked about a pose in which the arm was
\emph{physically against the mannequin}, one clearance model -- the one that
samples link \emph{origins} rather than link geometry, and which runs on the
physical arm during homing -- returned \SI{367.0}{\milli\metre}, the same
value it returns for a visibly clear pose. Every clearance figure in this
report computed through the guard's capsule model stands (the majority);
but the standing claim that wearer margin is enforced \emph{at all times}
was, at the time of writing, \textbf{not enforced in code on hardware}. This
is the one item in \cref{ch:conclusion}'s further-work list that is a safety
defect rather than an incompleteness, and it is placed above the mount
change and above the study for that reason.

\section{Three things the real cameras changed}

Beyond the headline numbers of \cref{ch:vision}: the surface-fitting stage
transferring to hardware within a factor of three of its simulated figure is
the one component of the perception chain now known to work outside
simulation, and the one the pick path leans on hardest. The scene camera is
a \emph{different} instrument from the wrist cameras, not merely a distant
one -- an eighth of the depth returns, a twentieth of the plane support --
and no autonomous behaviour should expect a support surface from it. And the
learned detector contributed nothing the classical colour-and-geometry path
did not; no result in this report rests on a learned detector, worth noting
given that a language-model-free, learned-detector-optional design was
argued for on cost grounds rather than accuracy grounds.

\section{Comparison against the planning report, in full}

Divergences fall into three groups. \emph{Superseded by measurement}: the
haptic objective (\cref{ch:discussionmain}). \emph{Deferred and still open}:
coordinated pick-and-place on the real platform -- the task set is specified
and verified geometrically and the runner exists, but the arms were not
available for the period in which it would have run. \emph{Added}: the
safety architecture, the degradation architecture, the workspace
characterisation and the pilot were not in the plan; the first two follow
directly from discovering that sensor faults dominate, the third from asking
whether the task set was reachable before running it.
